\section{Representaciones y aplicaciones adicionales}

Hasta el momento, gran parte de los trabajos con NCA se han centrado en la morfogénesis y en su capacidad para generar estructuras visuales autoorganizadas. Sin embargo, también han surgido enfoques que amplían el marco conceptual hacia nuevas formas de representación y aplicaciones en dominios específicos. A continuación, se presentan algunas de estas variantes.

Mattie Tesfaldet et al. \cite{ViTCA} extienden los NCA mediante mecanismos de atención inspirados en arquitecturas de \textit{transformers}, dando lugar a \textit{Vision Transformer Cellular Automata} (ViTCA). Esta arquitectura híbrida permite que las interacciones locales tengan acceso contextual a regiones más amplias a través de \textit{self-attention}, lo que habilita organización global sin incurrir en la complejidad cuadrática típica de la atención completa. ViTCA adapta el proceso de los ViT tradicionales, que dividen una imagen en parches y los convierten en secuencias de \textit{tokens}, al formato de una cuadrícula de celdas, donde cada celda almacena información visual, posicional y oculta. En lugar de aplicar atención global, restringe el cálculo a vecindarios locales, manteniendo así la naturaleza distribuida del sistema y reduciendo el costo computacional. El entrenamiento se realiza de forma iterativa, alternando entre inyectar entradas corruptas y actualizar parcialmente los estados celulares hasta reconstruir la imagen. De este modo, ViTCA combina eficiencia, capacidad de autoorganización y organización espacial coherente.

En un ámbito distinto, Jian Tang et al. \cite{JianNCA} emplean CNN para simular el desarrollo de microestructuras en la solidificación de metales. Cada celda del dominio contiene un estado de fase $P\in[0,1]$ (donde $p=0$ representa líquido y $p=1$ sólido), tres ángulos de Euler $\theta=(\alpha_1,\alpha_2,\alpha_3)$ que definen su orientación cristalográfica, y la magnitud de sub enfriamiento $\Delta T=T_{melting}-T$. La evolución se modela con un NCA que aplica recurrentemente una CNN para predecir los incrementos en fase, orientación y estados ocultos. Estos canales ocultos facilitan el aprendizaje de patrones temporales complejos sin necesidad de reglas explícitas. Con ello, el modelo reproduce fenómenos como crecimiento competitivo o direcciones preferentes, con una eficiencia computacional sustancialmente mayor que la de muchos CA tradicionales y con capacidad de generalización fuera del rango de entrenamiento.

Por otro lado, Gordon H. Y. Li et al. \cite{PNCA} proponen el modelo \textit{Photonic Neural Cellular Automata} (PNCA), que aprovecha la velocidad y la interconexión masiva de la fotónica para implementar NCA en hardware óptico. A diferencia de arquitecturas convencionales basadas en MLP o CNN, difíciles de trasladar a entornos fotónicos por sus conexiones densas, PNCA utiliza reglas locales aprendidas mediante técnicas de aprendizaje profundo. Esto permite realizar clasificación de imágenes en sistemas ópticos, donde la actualización celular se implementa mediante interferencia lineal de luz y activaciones ópticas no lineales. Los autores demuestran clasificación binaria con alta precisión utilizando un número muy reducido de parámetros fotónicos programables, mostrando además robustez frente a datos fuera de distribución. PNCA representa una vía prometedora para llevar el aprendizaje autoorganizado a plataformas fotónicas escalables.

En el ámbito de la biología computacional, Caitlin Grasso y Josh Bongard \cite{ENCA} introducen los \textit{Empowered Neural Cellular Automata} (ENCA), una extensión que incorpora el concepto de \textit{empowerment} —la cantidad de control que una célula puede ejercer sobre su entorno— formalizado mediante medidas de información mutua entre acciones pasadas y estados sensoriales futuros. Incluir \textit{empowerment} como objetivo adicional en tareas de morfogénesis mejora la aptitud evolutiva y produce patrones de crecimiento más robustos y adaptativos, posicionando a cada célula como un agente autónomo dentro de un sistema distribuido y abriendo vías para estudiar coordinación y comunicación celular desde una óptica computacional.

Finalmente, mientras que modelos como ViTCA, PNCA o ENCA amplían las capacidades arquitectónicas y las plataformas de aplicación de los NCA, Richardson et al. \cite{richardsonNCA} utilizan los NCA como herramienta para modelar fenómenos dinámicos complejos. Entrenando NCA mediante optimización basada en gradientes para reproducir trayectorias espacio-temporales generadas por ecuaciones diferenciales parciales (por ejemplo, Gray–Scott), este enfoque captura tanto estructuras transitorias como patrones estacionarios —incluyendo la formación de patrones tipo Turing— y aborda el problema inverso de inferir reglas locales a partir de comportamiento emergente. Con una complejidad de red mínima, los autores muestran capacidad de generalización, respeto por simetrías y robustez frente a perturbaciones, posicionando a los NCA como un marco explicativo y eficiente para sistemas biológicos y físicos.

